% Separate post-CP18 candidate. No unrestricted least-index conclusion.
\subsection{A finite-level neighborhood of each canonical CM lift}

\begin{lemma}[Equivariant continuity of formal torsion]
\label{lem:B-RD54-continuity}
\status{Candidate; complete proof written, independent audits pending.}
Let $L$ be complete discretely valued of mixed characteristic $(0,p)$,
with algebraically closed residue field $k$, and write $a=v_L(p)$.
Let $\mathcal F_u$ be a one-dimensional height-two universal formal
deformation over $W(k)[[u]]$, with $\mathcal F_0$ the canonical
unramified quadratic Lubin--Tate formal group. Fix an integer $n\ge1$.
If $x\in\mathfrak m_L$ satisfies
\[
                          v_L(x)>2na,                 \tag{B-RD54.1}
\]
then there is a canonical isomorphism of finite Galois modules
\[
               \mathcal F_0[p^n](\overline L)
                   \simeq\mathcal F_x[p^n](\overline L).
                                                        \tag{B-RD54.2}
\]
It matches roots whose coordinate difference has valuation greater
than $na$.

Suppose moreover that a finite group $\Gamma$ acts on $L$ and on
the universal family, that $u=0$ is fixed, and that the point $u=x$
is equivariant for these actions. The universal family action is
understood to include the isomorphisms of formal groups over the
induced automorphisms of the parameter ring. Then (B-RD54.2)
commutes with the induced semilinear descent actions as well.
\end{lemma}
\begin{proof}
Write $R=W(k)[[u]]$. Weierstrass preparation over the complete
local ring $R$ gives
\[
                  [p^n]_{\mathcal F_u}(T)=U_u(T)W_u(T),
\]
where $U_u(T)$ is an integral unit series and $W_u(T)$ is monic
of degree $p^{2n}$. All its nonleading coefficients are in
$(p,u)$. Uniqueness of preparation shows that its specialization
at $u=0$ is the prepared polynomial for $\mathcal F_0$; hence
every coefficient of $W_u-W_0$ is divisible by $u$.
The same divisibility holds for the difference of the two formal
group laws.

Let $E/L$ be the splitting field of $\mathcal F_0[p^n]$, and
extend $v_L$ to $E$ without changing its normalization. For
every torsion coordinate $\alpha$ in this field,
\[
                       v_L\bigl(W_0'(\alpha)\bigr)=na.
                                                        \tag{B-RD54.3}
\]
Indeed, if $w(T)\,dT$ is an invariant differential with $w(0)=1$,
then
\[
 w([p^n](T))[p^n]'(T)=p^n w(T).
\]
At a torsion point all factors except $p^n$ are units, and the
prepared unit $U_0(\alpha)$ is also a unit. This proves
(B-RD54.3), including at $\alpha=0$.

Put $b=v_L(x)$. Coefficientwise divisibility gives
\[
 v_L(W_x(\alpha))\ge b,
 \qquad v_L(W_x'(\alpha)-W_0'(\alpha))\ge b.
\]
Since $b>2na$, the derivative still has valuation $na$.
The elementary Newton--Hensel argument therefore gives a root
$\beta\in E$ of $W_x$ with
\[
                       v_L(\beta-\alpha)\ge b-na>na.
                                                        \tag{B-RD54.4}
\]
For completeness, after one correction the error valuation
is at least $2(b-na)$; subsequent corrections strictly improve
this bound. Completeness gives a root. In a ball whose difference
valuations exceed $na$, the linear term of the Taylor expansion
strictly dominates every higher term, so this root is unique.
This argument works with the extended rational-valued valuation
on $E$ and introduces no further field extension.

Distinct torsion points of $\mathcal F_0[p^n]$ have coordinate
differences of valuation at most $a/(p^2-1)$. To see this,
translate their difference using the integral formal group law;
its valuation is unchanged by translation. If its exact order
is $p^r$, the canonical Lubin--Tate equation gives valuation
\[
                      \frac{a}{(p^2-1)p^{2(r-1)}}.
\]
In particular this is at most $a/(p^2-1)<na$.
Thus the balls in (B-RD54.4) are disjoint. There are $p^{2n}$
of them, and they give all roots of the separable polynomial
$W_x$.

The matching respects addition. If $\alpha_1,\alpha_2$ match
$\beta_1,\beta_2$, integrality of the formal laws and the
coefficientwise difference divisible by $x$ show that
\[
 \mathcal F_x(\beta_1,\beta_2)
       -\mathcal F_0(\alpha_1,\alpha_2)
\]
has valuation greater than $na$. The first coordinate is again
a $p^n$-torsion point. Uniqueness in the matching ball proves
additivity, and hence (B-RD54.2). The same uniqueness proves
$G_L$-equivariance.

Finally let $g\in\Gamma$. In universal coordinates its formal
isomorphism specializes at $u=0$ to the corresponding automorphism
of $\mathcal F_0$. Their difference has every coefficient
divisible by $u$. Evaluating at the equivariant point $x$,
and using $v_L(gx)=v_L(x)$, shows that applying the two descent
maps to a matched pair again gives a pair with difference
valuation greater than $na$. Uniqueness proves compatibility
with descent. This uses the actual isomorphisms and their
cocycle, not merely the action on isomorphism classes. For the
universal height-two family these are supplied uniquely by
\cite[Propositions~14.7 and~14.9, including (14.8)--(14.10)]
{HopkinsGross1994Bundles}.
\end{proof}

\begin{proposition}[An excluded open neighborhood in the remaining test]
\label{prop:B-RD54-neighborhood}
\status{Candidate; complete proof written, independent audits pending.}
Use the normalized remaining $p^2$ comparison of
Theorem~\ref{thm:B-RD49}, with $p\ge5$ and $e_K=p-1$.
Thus the residue field is algebraically closed, $\mu_p\subset K$,
and the two potentially supersingular height-two groups have
the same nonsquare finite determinant character.

For each group choose a faithful good tame field $L_i$ and
identify its special formal group with the Honda formal group
so that its finite descent subgroup is contained in
$\mu_{p^2-1}\subset\mathbf Z_{p^2}^{\times}$.
Let $x_i\in\mathfrak m_{L_i}$ be its parameter in the universal
deformation based at the canonical unramified quadratic CM lift.
If
\[
                   v_{L_i}(x_i)>4e_{L_i}
                         \qquad(i=1,2),                \tag{B-RD54.5}
\]
where $v(0)=+\infty$, then isomorphic $p^2$-torsion forces
the two requested special characters to be equal.

Consequently every remaining different-character pair has,
for any such two choices of normalized framings, at least one
parameter satisfying
\[
                       0<v_{L_i}(x_i)\le4e_{L_i}.
                                                        \tag{B-RD54.6}
\]
In particular, its normalized deformation parameter cannot
lie arbitrarily deeply in the canonical CM neighborhood on
both sides of the comparison.
\end{proposition}
\begin{proof}
The faithful finite subgroups and integral conjugacy are those
of Lemma~\ref{lem:B-RD49-inertia} and the integral conjugacy
argument in Lemma~\ref{lem:B-RD52-CM}. The latter averaging
argument applies to a finite cyclic subgroup of the special
division order whether or not the given lift has CM.
The canonical lift is fixed by this subgroup. The universal
deformation action, with its unique marked isomorphisms, is
described in
\cite[Propositions~12.10, 14.7, 14.9 and~14.13]
{HopkinsGross1994Bundles}.

Tate full faithfulness extends the generic descent maps to the
actual good models, so their marked universal parameters are
equivariant. At $x_i=0$, the same descent datum gives a
potentially CM generic group $M_i^0/K$: its good formal group
is the canonical unramified quadratic lift, and its descent
maps are the corresponding commuting roots of unity. Generic
finite-level Galois descent is effective and compatible with
the transition maps, so this constructs a height-two
Barsotti--Tate group over $K$. Its finite special character is
the original one.

Lemma~\ref{lem:B-RD54-continuity}, with $n=2$, now gives
\[
                      M_i[p^2]\simeq M_i^0[p^2]
\]
over $K$, including the full descent action. An isomorphism
between the original two truncations therefore gives one
between the two potentially CM truncations. Apply
Theorem~\ref{thm:B-RD52}. This proves (B-RD54.5) and its
contrapositive (B-RD54.6).
\end{proof}

\begin{remark}[Scope of the neighborhood restriction]
\label{rem:B-RD54-scope}
\status{Candidate scope statement; the remaining finite interval is unresolved.}
The parameter is chosen after specifying the special Honda
framing and its finite descent subgroup; it is not an invariant
of the unframed formal group. The result is quantified for
any such fixed choices and does not identify different
coordinate neighborhoods without their framings.
It imposes an explicit upper bound on at least one nonzero
parameter in a hypothetical counterexample. It neither
constructs a counterexample within this interval nor proves
the unrestricted $p^2$ character criterion.
No Hasse-thickness estimate is transferred to an isogenous
quotient in this argument.
\end{remark}
